\documentclass[12pt,a4paper]{article}
\usepackage[UTF8]{ctex}
\usepackage{amsmath,amssymb,amsthm}
\usepackage{geometry}
\usepackage{bm}
\usepackage{hyperref}

\geometry{margin=2.5cm}

\title{例12.12\quad 销插入\quad 详解}
\author{现代机器人学：第12章 抓取与操作\\
详细解释}
\date{}

\begin{document}

\maketitle

\section{问题描述}

考虑力控制的平面销在插入过程中与孔的两点接触（如图12.28所示）。

\subsection{系统组成}

\begin{itemize}
    \item \textbf{销}：力控制的平面物体
    \item \textbf{孔}：固定约束
    \item \textbf{接触}：销与孔在两个点处接触
    \item \textbf{力控制}：力控制器对销施加特定的wrench
\end{itemize}

\section{问题设置}

\subsection{接触配置}

图12.28显示了：
\begin{itemize}
    \item 作用在销上的接触摩擦锥
    \item 相应的复合wrench锥
    \item 使用力矩标记说明
\end{itemize}

\subsection{力控制}

力控制器可以对销施加特定的wrench $F_{\text{control}}$。

\section{问题1：卡住现象}

\subsection{情况描述}

如果力控制器将wrench $F_1$施加到销上，它可能会\textbf{卡住}。

\subsection{卡住的机制}

\begin{itemize}
    \item \textbf{力控制器施加wrench} $F_1$
    \item \textbf{孔可能产生接触力}来平衡$F_1$
    \item \textbf{如果接触力在摩擦锥内}，则可以平衡$F_1$
    \item \textbf{结果}：销可能卡在这个位置，无法继续插入
\end{itemize}

\subsection{数学表述}

销卡住的条件是：存在接触wrench $\sum_i k_i F_i$，使得：

\begin{equation}
F_{\text{control}} + \sum_i k_i F_i = 0, \quad k_i \geq 0, F_i \in \mathcal{WC}_i
\end{equation}

即，接触可以产生与控制器wrench相反的wrench，从而平衡它。

\subsection{物理解释}

\begin{itemize}
    \item 控制器试图推动销插入
    \item 但接触力可以抵抗这个推动
    \item 如果接触力足够大，可以完全平衡控制力
    \item 销无法移动，卡在当前位置
\end{itemize}

\section{问题2：成功插入}

\subsection{情况描述}

如果力控制器将wrench $F_2$施加到销上，\textbf{插入继续进行}。

\subsection{成功插入的机制}

\begin{itemize}
    \item \textbf{力控制器施加wrench} $F_2$
    \item \textbf{接触无法平衡} $F_2$
    \item \textbf{结果}：销继续移动，插入继续进行
\end{itemize}

\subsection{数学表述}

插入成功的条件是：不存在接触wrench $\sum_i k_i F_i$，使得：

\begin{equation}
F_{\text{control}} + \sum_i k_i F_i = 0, \quad k_i \geq 0, F_i \in \mathcal{WC}_i
\end{equation}

即，接触无法产生与控制器wrench相反的wrench来平衡它。

\subsection{物理解释}

\begin{itemize}
    \item 控制器推动销插入
    \item 接触力无法完全抵抗这个推动
    \item 存在净wrench使销继续移动
    \item 插入成功
\end{itemize}

\section{关键区别}

\subsection{Wrench $F_1$ vs $F_2$}

\begin{itemize}
    \item \textbf{$F_1$}：在复合wrench锥内 $\Rightarrow$ 可能被平衡 $\Rightarrow$ 可能卡住
    \item \textbf{$F_2$}：不在复合wrench锥内 $\Rightarrow$ 无法被平衡 $\Rightarrow$ 插入继续
\end{itemize}

\subsection{力矩标记的视角}

使用力矩标记：
\begin{itemize}
    \item 如果控制wrench在复合wrench锥内，接触可以产生相反的wrench
    \item 如果控制wrench不在复合wrench锥内，接触无法产生相反的wrench
\end{itemize}

\section{力闭合与楔入}

\subsection{力闭合条件}

如果两个接触处的摩擦系数足够大，使得两个摩擦锥可以"看到"彼此的基（如图12.22(b)所示），则：

\begin{itemize}
    \item 销处于\textbf{力闭合}
    \item 接触可能能够抵抗任何wrench
    \item 这取决于两个接触之间的内力
\end{itemize}

\subsection{楔入现象}

当销处于力闭合时，它被称为\textbf{楔入}（wedged）。

\begin{itemize}
    \item \textbf{定义}：销被"卡"在孔中，无法移动
    \item \textbf{原因}：接触可以产生任何方向的wrench来抵抗外部wrench
    \item \textbf{条件}：两个摩擦锥可以"看到"彼此的基
\end{itemize}

\subsection{楔入的数学条件}

销楔入的条件是：
\begin{itemize}
    \item 两个接触的摩擦锥足够宽
    \item 两个摩擦锥可以"看到"彼此的基
    \item 四个接触wrench（每个接触两条摩擦锥边）正张成整个wrench空间
    \item 销处于力闭合
\end{itemize}

\section{插入策略}

\subsection{避免卡住}

为了避免卡住，应该：

\begin{enumerate}
    \item \textbf{选择合适的控制wrench}：
    \begin{itemize}
        \item 确保控制wrench不在复合wrench锥内
        \item 或者确保控制wrench的方向使接触无法平衡它
    \end{itemize}
    
    \item \textbf{控制插入速度}：
    \begin{itemize}
        \item 快速插入可能减少卡住的可能性
        \item 但需要平衡速度和精度
    \end{itemize}
    
    \item \textbf{使用适当的摩擦系数}：
    \begin{itemize}
        \item 较低的摩擦系数可以减少卡住的可能性
        \item 但可能影响其他操作性能
    \end{itemize}
\end{enumerate}

\subsection{成功插入的条件}

为了成功插入，需要：

\begin{itemize}
    \item 控制wrench不在复合wrench锥内
    \item 或者控制wrench的方向使接触无法完全平衡它
    \item 存在净wrench使销继续移动
\end{itemize}

\section{动态分析}

\subsection{准静态 vs 动态}

\begin{itemize}
    \item \textbf{准静态}：假设速度很小，忽略惯性力
    \item \textbf{动态}：考虑加速度和惯性力
\end{itemize}

\subsection{动态情况下的卡住}

在动态情况下：
\begin{itemize}
    \item 即使接触可以平衡控制wrench，销可能仍然有速度
    \item 惯性可能使销继续移动
    \item 但最终会减速并可能卡住
\end{itemize}

\section{实际应用}

\subsection{装配任务}

这个例子展示了：

\begin{itemize}
    \item \textbf{插入操作}：如何成功执行插入任务
    \item \textbf{避免卡住}：如何避免在插入过程中卡住
    \item \textbf{力控制策略}：如何选择合适的控制wrench
\end{itemize}

\subsection{相关应用}

\begin{itemize}
    \item \textbf{销插入}：将销插入孔中
    \item \textbf{螺栓拧入}：将螺栓拧入螺纹孔
    \item \textbf{装配操作}：各种装配任务
    \item \textbf{精密插入}：需要精确控制的操作
\end{itemize}

\section{关键要点总结}

\begin{enumerate}
    \item \textbf{卡住现象}：如果控制wrench在复合wrench锥内，销可能卡住
    
    \item \textbf{成功插入}：如果控制wrench不在复合wrench锥内，插入可以继续
    
    \item \textbf{力闭合}：如果两个摩擦锥可以"看到"彼此的基，销处于力闭合
    
    \item \textbf{楔入}：力闭合时，销可能被楔入，无法移动
    
    \item \textbf{插入策略}：选择合适的控制wrench和插入策略可以避免卡住
    
    \item \textbf{力矩标记}：使用力矩标记可以直观地理解wrench锥和插入过程
\end{enumerate}

\section{设计建议}

\subsection{避免卡住的设计}

\begin{itemize}
    \item \textbf{几何设计}：设计销和孔的几何形状，使接触wrench锥较小
    \item \textbf{材料选择}：选择适当的摩擦系数
    \item \textbf{控制策略}：选择合适的控制wrench方向
    \item \textbf{速度控制}：使用适当的速度来减少卡住的可能性
\end{itemize}

\subsection{成功插入的设计}

\begin{itemize}
    \item \textbf{引导}：使用引导结构帮助插入
    \item \textbf{对齐}：确保销和孔正确对齐
    \item \textbf{力控制}：使用力控制而不是位置控制
    \item \textbf{反馈}：使用力反馈来调整插入策略
\end{itemize}

\end{document}

